\chapter*{Vorwort}
\addcontentsline{toc}{chapter}{Vorwort}

\begin{merksatz}
\boxtitel{Die Zahlen dieser Reihe}
Fast jede Zahl der Bände 16 bis 26 stammt aus einer Simulation: Zwanzigtausend Irrfahrten, zweihunderttausend Versicherte, fünftausend Paare von Zeitreihen. Die Bände haben diese Zahlen benutzt, als wären sie berechnet. Dieser Band fragt, was sie wirklich sind.

\begin{center}
{\small
\renewcommand{\arraystretch}{1.5}
\begin{tabular}{@{}P{2.2cm}P{6.4cm}@{}}
\toprule
\textbf{Band 18} & Epidemien, gezogen aus Zufallszahlen, die ein Rechner erzeugt hat.\\
\textbf{Band 22} & Lernkurven über hundert Wiederholungen --- mit welchem Fehlerbalken?\\
\textbf{Band 25} & Brownsche Pfade mit $10^{6}$ Schritten --- wie nahe an der Theorie?\\
\textbf{Band 26} & Scheinregression in $77$ Prozent der Fälle --- plus oder minus wie viel?\\
\bottomrule
\end{tabular}}
\end{center}

\begin{center}
\begin{minipage}{0.9\textwidth}
\textbf{Eine Simulation ist ein Experiment. Sie hat Fehler, Verzerrungen und Grenzen wie jedes andere --- und sie verlangt dieselbe Sorgfalt.}
\end{minipage}
\end{center}
\end{merksatz}

\begin{merksatz}
\boxtitel{Drei Ergebnisse, die man nicht erwartet}
\begin{center}
{\small
\renewcommand{\arraystretch}{1.55}
\begin{tabular}{@{}P{1.0cm}P{7.4cm}@{}}
\toprule
$1$ & Ein Zufallsgenerator, der jahrzehntelang in Gebrauch war, erzeugt Punkte, die vollständig auf fünfzehn Ebenen liegen --- und bestand alle gängigen Tests. \kursiv{(Kapitel 2)}\\
$2$ & In zwanzig Dimensionen hat ein Gitter mit einer Million Punkten nur zwei Punkte je Achse; Zufallspunkte stören sich daran nicht. \kursiv{(Kapitel 3)}\\
$3$ & Man kann eine Simulation hundertmal genauer machen, ohne sie hundertmal länger laufen zu lassen --- indem man denselben Zufall zweimal benutzt. \kursiv{(Kapitel 5)}\\
\bottomrule
\end{tabular}}
\end{center}

\begin{center}
\begin{minipage}{0.88\textwidth}
\kursiv{Das zweite Ergebnis ist der Grund, warum es diesen Band gibt. \textbf{Oberhalb von vier oder fünf Dimensionen versagt jede systematische Rechnung}, und die Zufallsmethode wird vom Notbehelf zur einzigen Möglichkeit. Fast alle Fragen der zweiten Werkhälfte --- Integrale über Patientenmerkmale, über Pfade, über Parameter --- sind hochdimensional.}
\end{minipage}
\end{center}
\end{merksatz}

\begin{merksatz}
\boxtitel{Die sechs Teile}
\begin{center}
{\small
\renewcommand{\arraystretch}{1.5}
\begin{tabular}{@{}P{1.4cm}P{3.6cm}P{3.0cm}@{}}
\toprule
\textbf{Teil} & \textbf{Thema} & \textbf{Kapitel}\\
\midrule
I & Der Zufall im Rechner & $1$--$3$\\
II & Genauer werden & $4$--$6$\\
III & Ketten und Pfade & $7$--$9$\\
IV & Systeme über die Zeit & $10$--$12$\\
V & Die Praxis & $13$--$14$\\
VI & Stand & $15$\\
\bottomrule
\end{tabular}}
\end{center}

\schritt{Der Leitfall}
Eine Intensivstation mit zwölf Betten. Patienten treffen zufällig ein, bleiben unterschiedlich lange, und wer kein Bett findet, muss verlegt werden. Wie oft geschieht das --- und was ändert ein dreizehntes Bett?

\begin{center}
\begin{minipage}{0.88\textwidth}
\kursiv{\textbf{Die Frage ist einfach genug, um sie im Sonderfall exakt zu lösen, und realistisch genug, um das sofort zu verlieren}: Sobald Dringlichkeit, Tageszeit oder die Abhängigkeit der Verweildauer von der Auslastung hinzukommen, gibt es keine Formel mehr. Dann bleibt die Simulation --- und die Frage, wie weit man ihr trauen darf.}
\end{minipage}
\end{center}
\end{merksatz}

\begin{warnung}
\boxtitelr{Drei Warnungen}
\begin{center}
{\small
\renewcommand{\arraystretch}{1.55}
\begin{tabular}{@{}P{2.2cm}P{6.4cm}@{}}
\toprule
\textbf{Erste} & Eine Simulation ohne Fehlerbalken ist eine Anekdote. Jede Zahl braucht ihre Streuung.\\
\textbf{Zweite} & Der Fehler fällt mit $1/\sqrt n$. Zehnmal genauer heißt hundertmal länger --- es sei denn, man ist klug.\\
\textbf{Dritte} & Was man simuliert, ist das Modell. Die Simulation prüft die Rechnung, nicht die Annahmen.\\
\bottomrule
\end{tabular}}
\end{center}

\begin{center}
\begin{minipage}{0.88\textwidth}
\kursiv{Die dritte Warnung verbindet diesen Band mit Band 23 und 26. \textbf{Eine Simulation kann zeigen, dass ein Schätzer unter den eigenen Annahmen funktioniert} --- so wurde es in Band 26 durchgehend getan. Ob die Annahmen gelten, kann sie nicht zeigen; dafür gibt es keine Rechnung.}
\end{minipage}
\end{center}
\end{warnung}

\begin{merksatz}
\boxtitel{Was dieser Band voraussetzt}
\begin{center}
{\small
\renewcommand{\arraystretch}{1.48}
\begin{tabular}{@{}P{3.4cm}P{1.4cm}P{2.8cm}@{}}
\toprule
\textbf{Werkzeug} & \textbf{Band} & \textbf{gebraucht in}\\
\midrule
Wahrscheinlichkeit & $4$ & durchgehend\\
Numerik & $5$ & Kapitel 10--12\\
Statistik & $17$ & Kapitel 4, 13\\
Markow-Ketten & $25$ & Kapitel 7--9\\
Differential\-gleichungen & $14$ & Kapitel 10--11\\
\bottomrule
\end{tabular}}
\end{center}

\begin{center}
\begin{minipage}{0.88\textwidth}
\kursiv{\textbf{Alle Rechnungen dieses Bandes sind in wenigen Zeilen Python nachvollziehbar}; der Code steht im Anhang und im Quellenverzeichnis der Reihe. Das ist hier kein Zusatz, sondern der Gegenstand.}
\end{minipage}
\end{center}
\end{merksatz}
